%% ==================================================
%% Section: Circuits with measurements
%% Sources: arXiv:2307.01696 (Malz, Styliaris, Wei, Cirac), paragraphs
%% "Preparations using measurements" and "Long-range MPS using
%% measurements"; arXiv:2103.13367 (Piroli, Styliaris, Cirac), paragraphs
%% "Quantum circuits and LOCC" and "State transformations with QC and LOCC",
%% Definition "Transformations under QC and LOCC", and Example 1.
%% ==================================================

\section{Circuits with measurements}
\label{sec:qc_measurements}

The model of preparation with measurements is that of \cite{Piroli2021LOCC}:
a local circuit, measurements, and local unitaries depending on the outcomes
of all the measurements, which are communicated to every site at no cost.
Only the layers of the circuit are counted. In this model GHZ-type states are
prepared in constant depth, although their connected correlations do not
decay with distance, while by Lemma~\ref{lem:qc_light_cone_correlations} those
of a state prepared by a local circuit of depth $T$ vanish beyond distance
$2T$. Rounds of measurements are composed, and a vector prepared with
measurements may be followed by a local circuit that does not depend on the
outcomes.

\begin{definition}[Preparation with measurements]
    \label{def:qc_measurement_preparation}
    \lean{QuantumCircuit.outcomeProj, QuantumCircuit.outcomeProj_eq_ctrlProj,
        QuantumCircuit.MeasurementProtocol,
        QuantumCircuit.MeasurementProtocol.preMeasurement,
        QuantumCircuit.MeasurementProtocol.postMeasurement,
        QuantumCircuit.MeasurementProtocol.output,
        QuantumCircuit.MeasurementProtocol.IsPreparationOf,
        QuantumCircuit.IsPreparedWithMeasurementsInDepth}
    \leanok
    \uses{def:qc_local_circuit}
    The local labels are $\{0,\dots,d-1\}$. For a set $S$ of sites of the
    ring of $N$ sites and a string $m\colon S\to\{0,\dots,d-1\}$, let $P_m$
    be the projection onto the span of the basis vectors $\ket{x}$ with
    $x_i=m_i$ for every $i\in S$; these are the outcome projections of a
    measurement of the sites of $S$ in the computational basis, and $P_m$ is
    the projection onto the configurations agreeing on $S$ with any
    configuration extending $m$. A \emph{protocol} consists of a product
    vector $\ket{\pi}$, a local circuit $U$ of depth $T_1$, a set $S$ of
    measured sites, and for every $m$ unitaries $v_{m,i}$ on the sites $i$,
    with product $V_m=\bigotimes_iv_{m,i}$. It \emph{prepares} a vector
    $\ket{\psi}$ if $\ket{\pi}\neq0$ and for every $m$ with
    $P_mU\ket{\pi}\neq0$ there is $c\in\mathbb C$ with
    $V_mP_mU\ket{\pi}=c\ket{\psi}$. A vector is \emph{prepared with
        measurements in depth $T$} if a protocol with $T_1\leq T$ prepares
    it.

    The model is that of \cite{Piroli2021LOCC}, where the class of such
    preparations is called $\mathrm{QCcc}$. The correction $V_m$ may depend on
    the whole string $m$: classical communication is global and free, and the
    single-site corrections are not counted. Ancillas are sites of the ring,
    and the local unitaries between the layers of the circuit, acting on a
    site and its ancillas and free in \cite{Piroli2021LOCC}, are counted here
    as gates. The correction $V_m$ is a local unitary in the sense of
    \cite{Piroli2021LOCC}. So every preparation in this sense is one in the
    sense of \cite{Piroli2021LOCC}, of no larger depth.
\end{definition}

\begin{lemma}[Outcomes of nonzero probability]
    \label{lem:qc_measurement_outcomes}
    \lean{QuantumCircuit.outcomeProj_mulVec_apply,
        QuantumCircuit.sum_outcomeProj, Matrix.finKronecker_mem_unitary,
        Matrix.eq_zero_of_mem_unitary_of_mulVec_eq_zero,
        QuantumCircuit.MeasurementProtocol.sum_postMeasurement,
        QuantumCircuit.MeasurementProtocol.preMeasurement_ne_zero,
        QuantumCircuit.MeasurementProtocol.exists_postMeasurement_ne_zero,
        QuantumCircuit.MeasurementProtocol.output_ne_zero,
        QuantumCircuit.MeasurementProtocol.IsPreparationOf.ne_zero}
    \leanok
    \uses{def:qc_measurement_preparation}
    The outcome projections satisfy $\sum_mP_m=\mathbb 1$, and a product of
    unitaries on the sites is unitary; a unitary matrix maps only the zero
    vector to zero. Hence, for a protocol with
    $\ket{\pi}\neq0$, some outcome has $P_mU\ket{\pi}\neq0$, then
    $V_mP_mU\ket{\pi}\neq0$, and every vector that the protocol prepares is
    nonzero.
\end{lemma}
\begin{proof}
    \leanok
    For each configuration $x$ exactly one string $m$, the restriction of $x$
    to $S$, has $P_m\ket{x}\neq0$. So $\sum_mP_mU\ket{\pi}=U\ket{\pi}$, which
    is nonzero because $U$ is unitary; the unitary $V_m=\bigotimes_iv_{m,i}$,
    with $V_m^\dagger V_m=\bigotimes_iv_{m,i}^\dagger v_{m,i}=\mathbb 1$,
    keeps a nonzero vector nonzero, and $V_mP_mU\ket{\pi}=c\ket{\psi}\neq0$ forces
    $\ket{\psi}\neq0$.
\end{proof}

\begin{lemma}[Protocols without measurements]
    \label{lem:qc_measurement_unitary_case}
    \lean{QuantumCircuit.isPreparedWithMeasurementsInDepth_of_isPreparedInDepth,
        QuantumCircuit.exists_isPreparedInDepth_of_isPreparationOf_of_measured_eq_empty,
        QuantumCircuit.smul_productVector}
    \leanok
    \uses{def:qc_measurement_preparation, lem:qc_measurement_outcomes}
    A nonzero vector prepared in depth $T$ is prepared with measurements in
    depth $T$. Conversely, if a protocol measuring no site, with $T_1\leq T$,
    prepares $\ket{\psi}$, then there are unitaries $u_i$ on the sites $i$
    such that $\bigotimes_iu_i\ket{\psi}$ is prepared in depth $T$.
\end{lemma}
\begin{proof}
    \leanok
    For the first statement take $S=\emptyset$ and $v_{m,i}=\mathbb 1$. For
    the second, $P_m=\mathbb 1$ for the only string $m$, so
    $V_mU\ket{\pi}=c\ket{\psi}$ with $c\neq0$ by
    Lemma~\ref{lem:qc_measurement_outcomes}. Take $u_i=v_{m,i}^\dagger$. Then
    $\bigotimes_iu_i\ket{\psi}=U(c^{-1}\ket{\pi})$, where $c^{-1}\ket{\pi}$
    is the product vector with the vector of one site multiplied by $c^{-1}$,
    and the circuit $U$ of depth at most $T$ is padded with empty layers to
    depth $T$.
\end{proof}

\begin{lemma}[Layers of permutation gates]
    \label{lem:qc_permutation_layers}
    \lean{QuantumCircuit.IsLocalPerm,
        QuantumCircuit.IsLocalPerm.permMatrix_mem_supportedOperators,
        QuantumCircuit.shiftPerm, QuantumCircuit.shiftPerm_apply,
        QuantumCircuit.isLocalPerm_shiftPerm, QuantumCircuit.permLayer,
        QuantumCircuit.exists_permLayer_op_mulVec,
        QuantumCircuit.exists_shiftLayer_op_mulVec,
        Matrix.finKronecker_permMatrix_mulVec}
    \leanok
    \uses{def:qc_local_circuit}
    Let $\sigma$ be a permutation of the configurations of the ring that
    changes only the sites of a set $X$, the new values at these sites
    depending only on the old values there. Its permutation matrix $P_\sigma$,
    $(P_\sigma v)(x)=v(\sigma x)$, is a unitary acting on $X$. Examples are the
    shift $x\mapsto x+f(x)e_t$ at a site $t\in X$, with values in
    $\mathbb Z_d$ and $f$ reading only the sites of $X\setminus\{t\}$. If
    $\sigma_k$, $k\in K$, are such permutations for pairwise disjoint
    pairs $\{k,k+1\}$, the layer of the gates $P_{\sigma_k}$ acts as
    $v\mapsto v\circ\rho$, where $\rho(x)$ agrees with $\sigma_k(x)$ on the
    pair $\{k,k+1\}$ and with $x$ off all the pairs. For shifts
    $\sigma_k\colon x\mapsto x+f_k(x)e_{t_k}$ this reads
    $\rho(x)_{t_k}=x_{t_k}+f_k(x)$, and $\rho(x)_i=x_i$ at every site $i$
    that is none of the $t_k$. Finally, for permutations $\tau_i$ of the
    labels of the sites $i$, the product $\bigotimes_iP_{\tau_i}$ acts as
    $v\mapsto v\circ\tau$, where $\tau(x)_i=\tau_i(x_i)$.
\end{lemma}
\begin{proof}
    \leanok
    Enumerating $X$ writes $P_\sigma$ as $Y\otimes\mathbb 1$ with
    $Y_{u,w}=1$ if $w=\sigma(x)|_X$ for some $x$ with $x|_X=u$, and $0$
    otherwise; permutation matrices are unitary. For the layer, induct over
    the pairs: $P_{\sigma_a}P_\rho=P_{\rho\sigma_a}$, and $\rho\sigma_a$
    agrees with $\sigma_a$ on the pair $\{a,a+1\}$, which the earlier pairs
    do not meet, and with $\sigma_k$ on an earlier pair $\{k,k+1\}$, because
    $\sigma_a(x)$ agrees with $x$ there. The entry of
    $\bigotimes_iP_{\tau_i}$ at $(x,y)$ is $\prod_i[y_i=\tau_i(x_i)]$, which is
    $1$ if $y=\tau(x)$ and $0$ otherwise.
\end{proof}

From here on the local labels are $\mathbb Z_d=\{0,\dots,d-1\}$ with
$d\ge1$.

\begin{definition}[Measurement rounds]
    \label{def:qc_measurement_rounds}
    \lean{QuantumCircuit.MeasurementRound, QuantumCircuit.MeasurementRound.kraus,
        QuantumCircuit.MeasurementRound.depth,
        QuantumCircuit.MeasurementRound.toProtocol,
        QuantumCircuit.MeasurementRound.outputs,
        QuantumCircuit.IsPreparedWithMeasurementRoundsInDepth}
    \leanok
    \uses{def:qc_measurement_preparation}
    A \emph{measurement round} consists of a local circuit $U$, a set $S$ of
    sites measured in the computational basis, and for every outcome
    $m\colon S\to\mathbb Z_d$ unitaries $v_{m,i}$ on the sites $i$, with product
    $V_m=\bigotimes_{i}v_{m,i}$. Its operator for the outcome $m$ is $V_{m}P_{m}U$,
    with $P_m$ the outcome projection of
    Definition~\ref{def:qc_measurement_preparation}, and its depth is the
    number of layers of $U$. With a product vector it is a protocol in the sense
    of that definition. The \emph{outputs} of a sequence of rounds
    $R_1,\dots,R_n$ from a vector $\ket v$ are the vectors
    $K_n\cdots K_1\ket v$, with $K_j$ the operator of $R_j$ for some outcome. A
    vector $\ket\psi$ is \emph{prepared with measurement rounds in depth $T$}
    if there are a product vector $\ket\pi\neq0$ and rounds of total depth at
    most $T$ such that every nonzero output from $\ket\pi$ is a multiple of
    $\ket\psi$.

    Every round is a protocol of $\mathrm{QCcc}_\ell$ of
    \cite{Piroli2021LOCC} of no larger depth, in the sense of
    Definition~\ref{def:qc_measurement_preparation}; that paper performs a
    single round and notes that ``one could also define a more general scheme
    with multiple rounds of LOCC''. A sequence of rounds is a composition of
    such transformations, with the depth counted as the total number of layers.
    The number of rounds is not bounded: the tree circuit of
    Chapter~\ref{ch:log_depth_preparation} with $k$ levels uses $2k$ rounds,
    whereas in the class $\mathrm{QCcc}^{(k)}_\ell$ of the paragraph ``Phases
    of matter'' of \cite{Piroli2021LOCC} the number $k$ of composed
    transformations does not depend on the system size. Several rounds are needed when
    a correction must precede gates that do not commute with it, as for a
    teleported register on which an isometry acts.
\end{definition}

\begin{lemma}[Outputs of rounds]
    \label{lem:qc_measurement_rounds_outputs}
    \lean{QuantumCircuit.MeasurementRound.toProtocol_output,
        QuantumCircuit.MeasurementRound.exists_kraus_mulVec_ne_zero,
        QuantumCircuit.MeasurementRound.outputs_nil,
        QuantumCircuit.MeasurementRound.mem_outputs_cons,
        QuantumCircuit.MeasurementRound.mem_outputs_smul,
        QuantumCircuit.MeasurementRound.mem_outputs_append,
        QuantumCircuit.MeasurementRound.exists_mem_outputs_ne_zero,
        QuantumCircuit.IsPreparedWithMeasurementRoundsInDepth.ne_zero,
        QuantumCircuit.IsPreparedWithMeasurementRoundsInDepth.mono,
        QuantumCircuit.isPreparedWithMeasurementRoundsInDepth_of_isPreparedWithMeasurementsInDepth}
    \leanok
    \uses{def:qc_measurement_rounds, lem:qc_measurement_outcomes}
    The output of the protocol of a round and a product vector $\ket\pi$ for
    the outcome $m$ is $V_{m}P_{m}U\ket\pi$. For a nonzero vector some outcome of a
    round gives a nonzero vector, so a nonzero vector has a nonzero output
    under every sequence of rounds, and a vector prepared with measurement
    rounds is nonzero. The outputs of a multiple $c\ket v$ are the multiples by
    $c$ of the outputs of $\ket v$, and the outputs of a concatenation of two
    sequences are the outputs of the second sequence from the outputs of the
    first. A vector prepared with measurements in depth $T$ is prepared with
    one measurement round in depth $T$.
\end{lemma}
\begin{proof}
    \leanok
    Since $\sum_{m}P_{m}=\mathbb 1$ and $U$ is unitary, some $P_{m}U\ket v$ is
    nonzero, and the unitary $V_m$ keeps it nonzero. The other statements
    follow by induction over the sequence of rounds and from the linearity of
    the operators.
\end{proof}

\begin{definition}[Implementation by rounds]
    \label{def:qc_round_implementation}
    \lean{QuantumCircuit.MeasurementRound.IsImplementationOn,
        QuantumCircuit.MeasurementRound.IsRoundsImplementationOn}
    \leanok
    \uses{def:qc_measurement_rounds}
    Let $E$ be a set of vectors and $W$ a matrix. A round \emph{implements}
    $W$ on $E$ if for every outcome $m$ there is $c\in\mathbb C$ with
    $V_{m}P_{m}U\ket v=cW\ket v$ for every $\ket v\in E$. A sequence of rounds
    \emph{implements} $W$ on $E$ if every output from every $\ket v\in E$ is a
    multiple of $W\ket v$.
\end{definition}

\begin{lemma}[Implementations compose]
    \label{lem:qc_round_implementation_compose}
    \lean{QuantumCircuit.MeasurementRound.IsImplementationOn.exists_mem_outputs,
        QuantumCircuit.MeasurementRound.isRoundsImplementationOn_nil,
        QuantumCircuit.MeasurementRound.IsRoundsImplementationOn.mono,
        QuantumCircuit.MeasurementRound.IsRoundsImplementationOn.append,
        QuantumCircuit.MeasurementRound.IsRoundsImplementationOn.isPreparedWithMeasurementRoundsInDepth}
    \leanok
    \uses{def:qc_round_implementation, lem:qc_measurement_rounds_outputs}
    If a round $R$ implements $W$ on $E$ and $\ket v\in E$, every output of $R$
    followed by rounds $R'$ is a multiple of an output of $R'$ from $W\ket v$.
    The empty sequence implements $\mathbb 1$ on every set, and a sequence
    implementing $W$ on $E$ implements it on every subset of $E$. If a sequence
    $\mathcal R$ implements $W$ on $E$, $W$ maps $E$ into $E'$, and a sequence
    $\mathcal R'$ implements $W'$ on $E'$, then $\mathcal R$ followed by
    $\mathcal R'$ implements $W'W$ on $E$. If $\mathcal R$ implements $W$ on
    $E$ and the product vector $\ket\pi\neq0$ lies in $E$, then $W\ket\pi$ is
    prepared with measurement rounds in the total depth of $\mathcal R$.
\end{lemma}
\begin{proof}
    \leanok
    An output of the concatenation is an output $\ket w$ of $\mathcal R'$ from
    an output $c\,W\ket v$ of $\mathcal R$, hence $\ket w=c\ket{w'}$ with
    $\ket{w'}$ an output of $\mathcal R'$ from $W\ket v\in E'$, and
    $\ket{w'}=c'\,W'W\ket v$. The last statement is the definition of
    preparation with measurement rounds.
\end{proof}

\begin{definition}[Preparation with measurements and a circuit]
    \label{def:qc_measurement_circuit_preparation}
    \lean{QuantumCircuit.IsPreparedWithMeasurementsAndCircuitInDepth}
    \leanok
    \uses{def:qc_measurement_preparation, def:qc_local_circuit}
    A vector $\ket\psi$ is \emph{prepared with measurements and a circuit in
        depth $T$} if $\ket\psi=U\ket\varphi$ for a vector $\ket\varphi$
    prepared with measurements in depth $T_1$ and a local circuit $U$ of depth
    $T_2$, with $T_1+T_2\leq T$. The circuit $U$ does not depend on the
    outcomes of the measurements.
\end{definition}

\begin{lemma}[Layers acting on configurations]
    \label{lem:qc_configuration_layers}
    \lean{QuantumCircuit.apply_eq_ite_of_mulVec_eq_comp, QuantumCircuit.sigmaSite,
        QuantumCircuit.sigmaSite_injective, QuantumCircuit.layerCfg,
        QuantumCircuit.layerCfg_apply, QuantumCircuit.layerCfg_comp,
        QuantumCircuit.layerCfg_apply_of_forall_ne, QuantumCircuit.eq_layerCfg_iff,
        QuantumCircuit.noncommProd_embedOp_mulVec_eq_comp,
        QuantumCircuit.embedOp_mulVec_eq_comp, QuantumCircuit.permOp_mulVec,
        QuantumCircuit.embedOp_mulVec_productVector}
    \leanok
    \uses{def:qc_local_circuit, lem:qc_permutation_layers}
    Let $X_k$ be operators on pairwise disjoint sets $S_k$ of sites, each
    acting by a map $f_k$ of the configurations of $S_k$, that is
    $X_kv=v\circ f_k$ for every vector $v$; then $X_k$ has the entry $1$ at
    $(u,f_k(u))$ and $0$ elsewhere. The product of the $X_k$ acts by the map
    $\rho$ with $\rho(x)|_{S_k}=f_k(x|_{S_k})$ for every $k$ and
    $\rho(x)_i=x_i$ at every site $i$ outside all the $S_k$. A permutation
    $\tau$ of the sites acts as $v\mapsto v\circ(x\mapsto x\circ\tau)$. If $X$
    acts on the sites of $S$ and the product vector $\ket\pi$ is $\ket z$ on
    every site of $S$, then $(X\ket\pi)(y)=X_{y|_S,(z,\dots,z)}\pi'(y)$, where
    $\ket{\pi'}$ is the product vector with the all-ones vector at the sites
    of $S$ and the vectors of $\ket\pi$ elsewhere.
\end{lemma}
\begin{proof}
    \leanok
    Applying $X_k$ to the basis vector $\ket u$ gives the entries. The entry of
    the product at $(x,y)$ is $\prod_k[y|_{S_k}=f_k(x|_{S_k})]$ when $x$ and
    $y$ agree outside the $S_k$, and $0$ otherwise; this is $1$ exactly when
    $y=\rho(x)$. For the product vector, only the configuration $y'$ agreeing
    with $y$ outside $S$ and equal to $(z,\dots,z)$ on $S$ contributes.
\end{proof}

\begin{definition}[GHZ-type states]
    \label{def:qc_ghz_state}
    \lean{QuantumCircuit.ghzState, QuantumCircuit.ghzState_apply}
    \leanok
    Let $b,M\geq0$ and $\alpha\in\mathbb C^b$. The \emph{GHZ-type
        state} on $M$ sites of dimension $b$ is
    $\ket{\chi_M}=\sum_{i=0}^{b-1}\alpha_i\ket{i}^{\otimes M}$. For $M\geq1$
    and the sites numbered $0,\dots,M-1$, its amplitude at a configuration $s$
    is $\alpha_{s_0}$ if all $s_n$ are equal and $0$ otherwise.
\end{definition}

\begin{theorem}[GHZ-type states in constant depth with measurements]
    \label{thm:qc_ghz_measurement}
    \lean{QuantumCircuit.ghzSite, QuantumCircuit.withZeroAncillas,
        QuantumCircuit.ghzInitial, QuantumCircuit.ghzProtocol,
        QuantumCircuit.ghzProtocol_output,
        QuantumCircuit.isPreparedWithMeasurementsInDepth_withZeroAncillas_ghzState}
    \leanok
    \uses{def:qc_measurement_preparation, lem:qc_permutation_layers,
        def:qc_ghz_state}
    Let $M\geq1$, $b\geq1$ and $\alpha\in\mathbb C^b$ nonzero, with the
    labels $\{0,\dots,b-1\}=\mathbb Z_b$. On the ring of $2M$ sites of
    dimension $b$, numbered $0,\dots,2M-1$, place the system qudit
    $n\in\{0,\dots,M-1\}$ at the site $2n$ and its ancilla at the site
    $2n+1$. The GHZ-type state
    $\ket{\chi_M}=\sum_{i\in\mathbb Z_b}\alpha_i\ket{i}^{\otimes M}$ of
    Definition~\ref{def:qc_ghz_state} on the system qudits, with every ancilla in $\ket{0}$, is prepared with
    measurements in depth $2$. In the protocol below every outcome gives
    exactly this state. This is the construction of
    \cite[Example~1]{Piroli2021LOCC}, a preparation by $\mathrm{QCcc}_2$, for
    qudits, and the preparation of $\ket{\chi_M}$ in constant depth stated in
    the
    paragraph ``Long-range MPS using measurements'' of
    \cite{Malz2024Preparation}.
\end{theorem}
\begin{proof}
    \leanok
    The amplitude of $\ket{\chi_M}$ at $s$ is $\alpha_{s_0}$ if all $s_n$ are
    equal and $0$ otherwise. Values are taken in $\mathbb Z_b$. Start from the
    product vector with $\sum_i\alpha_i\ket{i}$ on the system qudit $0$,
    $\sum_i\ket{i}$ on the other system qudits, and $\ket{0}$ on the ancillas.
    A first layer of gates $\ket{k}\ket{a}\mapsto\ket{k}\ket{a+k}$ on the
    system qudit $n$ and its ancilla, and a second layer of gates
    $\ket{a}\ket{k}\mapsto\ket{a-k}\ket{k}$ on the ancilla $n$ and the system
    qudit $n+1$, for $n<M-1$, give
    \begin{align*}
        \sum_{k}\alpha_{k_0}\bigotimes_{n=0}^{M-1}\ket{k_n}
        \bigotimes_{n=0}^{M-2}\ket{k_n-k_{n+1}}\otimes\ket{0},
    \end{align*}
    the last factor being the last ancilla. Measure the ancillas $n<M-1$. The
    outcome $m$ forces $k_{n+1}=k_n-m_n$, so the system qudits are left in
    $\sum_i\alpha_i\bigotimes_n\ket{i-p_n}$ with the prefix sums
    $p_n=m_0+\cdots+m_{n-1}$. The single-site corrections, $X^{p_n}$ on the
    system qudit $n$ and $X^{-m_n}$ on the ancilla $n$, give $\ket{\chi_M}$
    with every ancilla in $\ket{0}$. The two layers are layers of permutation
    gates and the corrections are single-site permutations, and
    Lemma~\ref{lem:qc_permutation_layers} computes their action
    configuration by configuration. The product vector is nonzero because
    $\alpha\neq0$.
\end{proof}
